\Needspace{15\baselineskip}
\section{The Four Principal Eucharistic Prayers}

The 2008 Missal does not print four interchangeable lengths of one composition. It receives the Roman Canon and three prayers promulgated in 1968, each with its own order, historical relation, diction, and rules of selection. GIRM 365 recommends without ranking their sacramental efficacy: I may always be used and is especially apt in stated circumstances; II is especially suited to weekdays and particular circumstances; III is particularly apt for Sundays and festive days; IV carries an invariable Preface and can be used only where another Preface is not required. A lawful selection receives a complete prayer. One may not move a favored paragraph, saint list, epiclesis, or intercession from one into another. The overviews below classify architecture, rhetoric, and doctrine at a deliberate distance; they do not reproduce or sequentially translate the received clauses.

\ordermovement{\latin{Prex eucharistica I seu Canon Romanus}}{The Received Roman Canon: Concrete Names around the One Victim}{Approved principal option; variable Preface, proper inserts and commemorations, with selection guidance}{GIRM 365a; \latin{Ordo Missae}, Eucharistic Prayer I}

After the selected Preface and common Sanctus, \latin{Te igitur} opens a succession of petitions linked more by sustained address and sacrificial reference than by narrative transitions. Gifts are commended; the Church is remembered in visible communion; living persons and the assembly are placed before God; Mary, Joseph, apostles, martyrs, and all the saints are named; the household's offering is commended; and God is asked to act upon the oblation. The institution narrative forms the center. Anamnesis and offering then lead through the accepted offerings of Abel, Abraham, and Melchizedek to the heavenly-altar petition, remembrance of the dead, the sinners' hope of fellowship with the saints, blessing of creation through Christ, and the Trinitarian doxology.

This accumulation is its rhetoric. Repeated pleas for gracious acceptance do not imply that Christ might be an inadequate Victim or that the Father must be persuaded to love. They place gifts, offerers, ministers, local Church, dead, and sinners beneath the one Son's acceptable self-offering. The Canon does not narrate salvation history at length because it prays from within the economy already named by Preface, institution, memorial, scriptural types, saints, and final mediation.

\paragraph{A Roman prayer with several historical strata.}
Paul VI's \emph{Missale Romanum} gives the suitably broad official judgment that the Roman Canon had substantially assumed its form during the fourth and fifth centuries. The evidence does not date every clause that early. \emph{De sacramentis} IV.5.21--24 and IV.6.26--27 securely preserves a close late-fourth-century Western sequence beginning in the region of \latin{Quam oblationem}: pre-institution petition, institution recital, anamnesis and offering, the three scriptural exemplars, and the heavenly-altar plea. It does not transmit the entire opening, both saint catalogues, or every late conclusion in the received order. The early medieval sacramentaries provide secure complete or nearly complete Roman evidence only after further development. No single fourth-century author of the full Canon can be named.

The prayer's long predominance in the Roman rite likewise arose through reception rather than one decree composing it. Medieval additions and ceremonial developments surrounded a textual core that became highly stable. The postconciliar form remains recognizably the Canon while making real changes: the institution narratives of all four prayers were conformed; \latin{mysterium fidei} left the chalice formula and introduced the people's acclamation; several conclusions and parts of the saint lists became optional; and the renewed Order redistributed audible participation and ceremonial. Saint Joseph had already entered \latin{Communicantes} by John XXIII's decree of 13 November 1962. Historical continuity should neither conceal these revisions nor be dissolved by them.

\paragraph{The Church becomes a set of names.}
At \latin{Te igitur} and the first \latin{Memento}, the universal Church is not imagined apart from pope, local bishop, concrete worshippers, and those for whom prayer is offered. First Timothy~2:1--6 binds broad intercession to the one mediator; Hebrews~12:18--24 joins the pilgrim assembly to angels, the perfected righteous, and Jesus. The Canon gives those relations verbal density. Its names do not create rival mediations. Christ alone makes possible the communion in which saints are remembered and sinners ask to share.

The paired catalogues have a Roman accent: apostles and martyrs, women and men, figures universal in cult and figures closely bound to the city's memory. They are not an exhaustive census of heaven. The optional brackets of the current text should not be mistaken for permission to replace them with private lists. Proper forms of \latin{Communicantes} and \latin{Hanc igitur} allow Christmas, Epiphany, the Paschal celebrations, Ascension, Pentecost, and specified ritual occasions to enter the prayer according to the Missal rather than improvisation.

\paragraph{Household, oblation, and consecration.}
\latin{Hanc igitur} commends the offering of the Church's service and the whole household; the Priest extends his hands over the offerings. \latin{Quam oblationem} sharpens the petition toward the gifts themselves, asking divine action so that they become Christ's Body and Blood. The two passages should not be merged. One places an offering-bearing people under mercy; the other is a functional epiclesis over the oblation. It does not explicitly name the Holy Spirit and must not be made to recite the later prayers' pneumatic syntax.

Ambrose's close predecessor to \latin{Quam oblationem} helps interpret the sequence. Divine petition and Christ's effective words belong together. Aquinas later reads the Canon in the same ordered unity: the Priest speaks Christ's words in Christ's person, while the petitions dispose gifts and Church toward sacramental fruit (\emph{ST} III, q.~78; q.~83, a.~4). The Canon is not ``without an epiclesis'' if epiclesis names its consecratory petition, but it is without an explicit invocation of the Spirit over the gifts. Both statements are needed.

\paragraph{Accepted sacrifices and the heavenly altar.}
\latin{Unde et memores} remembers Passion, Resurrection, and Ascension and offers the life-giving Victim from gifts first received. \latin{Supra quae} recalls Abel's gift, Abraham's sacrifice, and Melchizedek's bread-and-wine offering (Gen.~4:4; 14:18--20; 22:1--18; Heb.~7). These figures typologically illuminate just offering, obedient faith, royal priesthood, and divine blessing. They are not previous Eucharistic consecrations, nor are they standards superior to Christ. Aquinas therefore locates the comparison in the offerers' acceptance rather than in any greater value of their gifts (\emph{ST} III, q.~83, a.~4, ad~8).

At \latin{Supplices te rogamus} the Priest bows, and the earthly altar opens verbally toward the heavenly. Revelation~8:3--4 gives a biblical field of angel, incense, prayer, and heavenly altar. The Canon does not identify its singular angel beyond dispute; angelic and Christological readings have both been received. Nor does the petition imagine Christ moving spatially from an inadequate earthly location to an acceptable heavenly one. Its stated fruit is that those who receive from this altar be filled with heavenly blessing and grace. Sacrifice tends toward Communion.

The second \latin{Memento} crosses death in resurrection hope; \latin{Nobis quoque peccatoribus} places the Priest and people among those who depend upon mercy rather than merit. \latin{Per quem haec omnia} restores the created horizon: through the Word by whom all things were made, created gifts are sanctified, enlivened, blessed, and bestowed. The doxology reveals the arc. The Canon has moved through particular names and gifts so that all may return to the Father through, with, and in the Son, in the Spirit's unity.

\begin{massbox}{The Canon's Roman abundance}
What can sound repetitive is a theology of inclusion. Gift, household, hierarchy, living, saints, dead, sinners, biblical offerers, earthly altar, heavenly blessing, and creation are successively drawn around Christ's words. The prayer does not escape the concrete in order to reach the universal; it reaches catholic breadth by naming particulars within the one sacrifice.
\end{massbox}

\ordermovement{\latin{Prex eucharistica II}}{Ancient Anaphoral Material Freely Recast in a Concise Roman Prayer}{Approved principal option; especially suited to weekdays and special circumstances; its own Preface or another qualifying Preface; proper deceased insertion}{GIRM 365b; \latin{Ordo Missae}, Eucharistic Prayer II}

Eucharistic Prayer II is the shortest principal prayer, but brevity is not theological emptiness. Its own Preface gives the compact anaphora a christological and Paschal horizon. After the Sanctus, divine holiness, explicit pneumatic consecratory language, received institution, ministerial thanksgiving, ecclesial communion, intercession, and eschatological hope are held within a notably linear design. These are formal and theological coordinates, not a substitute version of the prayer; the common anaphoral actions themselves were analyzed in the preceding chapter.

Its own Preface belongs especially naturally to the prayer, but GIRM 365b permits other Prefaces, particularly those presenting a concise synthesis of salvation. The permission is wider than a fixed-Preface rule and narrower than an assumption that any attractive Preface can be attached. The prayer's recommendation for weekdays and special circumstances is pastoral guidance, not a prohibition against every Sunday use. Concision may suit a celebration; convenience alone should not become the theology of selection.

\paragraph{What comes from the so-called \emph{Apostolic Tradition}?}
The modern prayer takes important language and its basic thanksgiving--institution--anamnesis/oblation--epiclesis--doxology sequence from the anaphoral material in chapter~4 of the church order conventionally called the \emph{Apostolic Tradition}. Twentieth-century scholarship commonly associated that work with Hippolytus of Rome and dated it near the beginning of the third century. Those claims can no longer be repeated as settled. The extant work survives through later translations and derivative church orders; its redaction is composite, its attribution to Hippolytus disputed, its Roman localization uncertain, and its material need not all come from a single early-third-century liturgy.

Nor did the 1967 study group simply translate chapter~4. The ancient witness lacks the present Roman Sanctus and post-Sanctus arrangement, the divided explicit epiclesis, the standardized institution narrative, the people's memorial acclamation, and the developed intercessions. Material was moved, supplemented, and reshaped. Pierre Jounel served as relator; Bernard Botte and Louis Bouyer made significant contributions; the promulgated prayer was nevertheless a corporate Coetus X redaction reviewed and approved within the Consilium and by Paul VI\@. Anecdotes that reduce its composition to two men improvising in haste cannot substitute for the redactional record.

The historically exact formula is therefore: a twentieth-century Roman Eucharistic Prayer freely adapting an ancient, textually complex anaphoral witness. It is neither an untouched ``Canon of Hippolytus'' nor a composition with no traditional roots.

\paragraph{A compressed Paschal economy.}
The own Preface's movement through Word, Incarnation, Passion, and Resurrection gives the prayer its narrative horizon before the Sanctus. Its brevity after the Sanctus depends upon that prior thanksgiving. When another admissible Preface is used, the prayer receives that day's economy rather than silently preserving every theme of its own.

The explicit epiclesis invokes the Spirit upon the gifts with the biblical image of dew. Dew can signify unobtrusive divine fecundity and eschatological life (Hos.~14:5; Isa.~26:19), but the prayer is not a code in which one natural symbol explains consecration. The subject is the Father sanctifying by the Spirit so that the gifts become the Son's Body and Blood. Christ's narrative then discloses what the petition seeks according to his command.

After the memorial acclamation, the prayer's thanksgiving for standing in divine presence retains a notable feature of its ancient source. Biblical standing can signify service and liturgical readiness rather than casual ease (Deut.~10:8; Rev.~7:9--15). Those admitted to minister do not congratulate themselves; they give thanks for a vocation received. The offering is the Paschal Christ, and the communion petition turns reception outward into the Spirit's unity.

The short intercessions then make visible what the prayer's center already implies. The local assembly is not the whole Church; the present generation is not the whole communion; the dead are not outside Christ's Pasch; Mary and the saints are not spectators. The petition for admission into their fellowship closes a concise prayer with eschatological amplitude.

\begin{studybox}{Brevity and completeness}
Eucharistic Prayer II contains all the constitutive actions named by GIRM 79. It should not be defended by pretending that every modern clause is ancient, nor dismissed because its verbal proportions differ from the Roman Canon. Its received theological form is ancient material transformed through a modern Roman composition.
\end{studybox}

\ordermovement{\latin{Prex eucharistica III}}{The Spirit Gathers a Worldwide Offering into Christ}{Approved principal option; may be used with any qualifying Preface and is especially suited to Sundays and festive days; proper deceased insertion}{GIRM 365c; \latin{Ordo Missae}, Eucharistic Prayer III}

Eucharistic Prayer III was composed for the reformed Missal rather than recovered as a complete ancient text. Cipriano Vagaggini prepared the foundational draft; Coetus X significantly revised it, and Paul VI reviewed the new prayers. Its idiom is Roman and biblical, while its broad sequence and pneumatology converse with Eastern anaphoral traditions. Naming Vagaggini's formative role is warranted; calling him sole author of every received sentence is not.

Because any qualifying Preface may precede it, Eucharistic Prayer III does not depend upon one fixed preliminary narration. Its distinctive grammar relates the Father's sanctifying action in creation to the gathering of a catholic people and a worldwide pure offering. The two pneumatic poles are especially important: sanctification of the gifts and transformation of communicants belong to one economy without becoming the same petition. Its intercessory breadth then tests the local assembly against Church, world, dead, saints, and Kingdom. This description identifies relations among elements rather than rendering their received wording.

\paragraph{From the nations to one pure offering.}
Malachi~1:11 judges corrupt priestly worship while announcing divine greatness among the nations and a pure offering everywhere. The \emph{Didache} receives the prophecy in its Eucharistic discipline (14.1--3); Justin reads it against pagan sacrifice and in relation to Christian thanksgiving (\emph{Dialogue with Trypho} 41); Irenaeus joins it to the new covenant's oblation offered throughout the world (\emph{Against Heresies} IV.17.5--18.4). Their interpretations are not identical and do not prove that Malachi predicted the wording of Eucharistic Prayer III\@. They give ancient warrant for its governing conjunction: catholic mission and pure worship converge in Christ's one offering.

The prayer first confesses that the Father gathers a people; only then does the local assembly ask for the gifts' sanctification. Catholicity is divine action before it is demographic reach. The offering from east to west is not millions of independent sacrifices competing with Calvary. In many times and places the Church sacramentally offers the one Victim through the one priesthood of Christ.

\paragraph{Two invocations, one transforming Spirit.}
Before the institution narrative the prayer asks the Spirit to sanctify the gifts. After anamnesis and oblation it asks the Father to look upon the Church's offering, recognize the Victim by whose sacrifice reconciliation was accomplished, and grant that those nourished by Christ be filled with the Spirit. Ephesians~4:4 supplies the compact goal: one body and one Spirit. Romans~12:1 supplies the moral horizon of persons offered as a living sacrifice. The second petition therefore does not reconsecrate the gifts. It asks that sacramental Communion achieve ecclesial transformation.

The next movement intensifies the claim: the Spirit is asked to make the communicants an enduring offering so that they receive the inheritance with the saints. Augustine's account of the whole redeemed city becoming sacrifice in Christ helps guard the meaning (\emph{City of God} X.6, X.20). Christians are not added to the altar as rival victims, nor does moral effort complete an insufficient Cross. Incorporated into the Son, they become by grace what the sacrament signifies.

\paragraph{Reconciliation opens into Kingdom.}
The prayer's offering is explicitly the Victim by whose sacrifice the Father reconciles the world. Second Corinthians~5:18--21 locates reconciliation in God's action through Christ and immediately entrusts an apostolic ministry of reconciliation. Colossians~1:15--20 extends peace through the blood of the Cross toward all creation. Eucharistic Prayer III similarly moves from sacrifice to the Church's peace and salvation, the world, pilgrims, departed, and final inheritance.

Its intercession does not turn the Kingdom into a vague happy ending. Mary, Joseph, apostles, martyrs, saints, and the departed entrusted to divine mercy locate hope in sanctified persons and God's judgment. The proper insertion for a deceased person can make suffrage concrete without changing the prayer into a private requiem. The closing vision gathers dispersed children into the inheritance where creation and Church reach their Christological end.

\begin{massbox}{A gathered people becomes an offering}
The distinctive arc is pneumatic and catholic: the Spirit sanctifies creation, gathers worshippers from every place, consecrates gifts, fills communicants, forms one body, and makes the Church an enduring offering in Christ. Mission, sacrifice, Communion, and eschatology are one movement rather than successive themes.
\end{massbox}

\ordermovement{\latin{Prex eucharistica IV}}{Creation and Covenant Carried through Christ into the New Creation}{Approved principal option with an inseparable fixed Preface; excluded where another Preface is required; no special deceased insertion}{GIRM 365d; \latin{Ordo Missae}, Eucharistic Prayer IV}

Eucharistic Prayer IV is the principal prayer whose Preface cannot be exchanged. Preface, Sanctus, and post-Sanctus form one continuous history of salvation. It is therefore available when the Mass has no Preface of its own and particularly on Sundays in Ordinary Time, not as a way to displace a required seasonal or proper thanksgiving. Its extended structure also explains why the Missal does not supply the special formula for a particular deceased person admitted in II and III.

The fixed Preface and post-Sanctus form one theological argument rather than a detachable seasonal opening. Creation as gift, the tragedy of sin, divine fidelity through covenant and prophecy, the incarnate Son's Pasch, and the Spirit's mission are narrated as one economy before the prayer concentrates that economy sacramentally upon the gifts. This broad map is needed to explain the fixed-Preface rule; it deliberately avoids a clause-by-clause surrogate for the approved text.

\paragraph{An Eastern-inspired composition, not a recovered Basilian text.}
The Consilium initially considered ancient Eastern models while seeking a prayer with a full history of salvation and an integral Preface. A subcommission, with Joseph Gelineau and Cipriano Vagaggini as important figures, produced a new composition. Alexandrian Basil was a significant comparative source; the broad narration of the economy has Antiochene affinities; the received Roman arrangement retains an explicit pre-institution epiclesis and post-institution communion petition. No single ancient anaphora is simply translated. ``Basil-influenced, newly composed'' is supportable; ``the Anaphora of St Basil restored to Rome'' is not.

This source history matters doctrinally. Patristic resonances belong to a modern synthesis approved by Paul VI\@. They are not archaeological guarantees that every sentence was prayed in fourth-century Alexandria or Antioch. Conversely, modern composition does not make its theology private. Promulgation receives the prayer into the Roman Church's lex orandi.

\paragraph{Creation remains inside redemption.}
Genesis~1:26--31 supplies image, stewardship, and created goodness; Genesis~3 supplies disobedience and exile. The prayer refuses two reductions. Sin does not reveal matter as evil, and creation's goodness does not make redemption unnecessary. Humanity's vocation is cosmic and priestly: rational creatures are to govern in obedience and return praise. When the image-bearer turns away, God continues to seek and teach.

Irenaeus's account of divine pedagogy illuminates the long sequence. God accustoms humanity to communion through covenants, prophets, and the Son, while the Word recapitulates human life in obedience (\emph{Against Heresies} III.18--22; IV.14--16). The prayer does not reproduce Irenaeus's argument, but both resist a salvation imagined as divine abandonment of the first creation. Fulfillment heals and elevates what the Creator intended.

The Christological center echoes Luke~4:16--21 and Isaiah~61: good news reaches poor, captivity, sorrow, and oppression. Hebrews~4:15 guards full solidarity without sin. Philippians~2:6--11 gives the descending obedience and exaltation by which lordship is revealed through self-emptying. The prayer's history is therefore not a catalogue of past interventions. It identifies the crucified Son as the truth of divine power and the pattern of ecclesial mission.

\paragraph{The Spirit completes Christ's work.}
The Spirit's mission appears twice before the institution narrative: first in the Incarnation, then in the sanctification of those who no longer live for themselves. Romans~14:7--9 and 2~Corinthians~5:15 give the Paschal logic of living for the risen Lord; John~14--16 gives the Spirit's mission of remembrance, witness, and guidance. The subsequent epiclesis is thus not a sudden ritual mechanism. It is the sacramental concentration of the same economy: the Father who sent the Spirit upon the Virgin and the Church is asked to sanctify the gifts.

The prayer's anamnetic field is unusually ample, extending the Paschal horizon through Christ's descent among the dead, Resurrection, Ascension, and awaited coming. The descent confesses the reach of the Pasch into the realm of death; it does not stage another saving event apart from the Cross. The oblation is correspondingly universal in horizon because the offered Christ's sacrifice is saving for the world.

Its communion and intercessory fields relate Spirit-given ecclesial unity to a scope wider than the persons visibly present. Ordered ministry and concrete assembly remain named realities, yet the prayer also looks beyond readily mapped ecclesial boundaries. That breadth neither declares every religion equivalent nor canonizes undefined sincerity. It entrusts persons to the God who searches hearts and whose universal salvific will is accomplished only through Christ.

The eschatological horizon likewise joins humble remembrance of the dead, communion with the saints, and creation's liberation. Romans~8:18--25 gives the canonical field: creation groans for the freedom of God's children. Salvation history reaches its goal not when creatures disappear, but when renewed creation becomes praise.

\begin{studybox}{Why the Preface cannot be exchanged}
The prayer's first movement names the transcendent Creator and angelic praise; its post-Sanctus then begins with humanity and follows one economy through fall, covenant, Incarnation, Pasch, and Spirit. Removing that Preface would remove the opening movement of the narrative, not merely exchange seasonal color.
\end{studybox}

The four prayers thus manifest one Eucharistic mystery through distinct grammars. The Canon gathers concrete persons and petitions around Christ's words; II compresses ancient anaphoral material into a concise modern Roman form; III makes the Spirit's worldwide gathering and the Church's transformed offering especially explicit; IV narrates creation and covenant into Pasch and new creation. Their plurality is received order, not permission to construct a fifth principal prayer from preferred parts.
